% Fonts and languages (LuaLaTeX + fontspec).
%
% This book is written in Vietnamese. Vietnamese is the main language, so babel
% supplies the section names (Chương, Mục lục, Phụ lục, Hình, Bảng, Chỉ mục)
% and the hyphenation patterns from hyph-vi. English stays loaded as a
% secondary language: every translated term carries its English original in
% parentheses, and quotations from the 32 corpus papers need English
% hyphenation to break.
%
% Fonts are loaded by FILENAME (texgyre*.otf in the TeX tree), not by display
% name - filename loading works on every machine without depending on the
% luaotfload font-name database being up to date. The TeX Gyre fonts already
% cover the whole Vietnamese repertoire, including the stacked tone-plus-vowel
% combinations (u+horn+dot-below and friends) that thinner fonts drop.

\usepackage[english]{babel}
\usepackage{fontspec}

\setmainfont{texgyrepagella}[
  Extension      = .otf,
  UprightFont    = *-regular,
  BoldFont       = *-bold,
  ItalicFont     = *-italic,
  BoldItalicFont = *-bolditalic,
]
\setsansfont{texgyreheros}[
  Extension      = .otf,
  UprightFont    = *-regular,
  BoldFont       = *-bold,
  ItalicFont     = *-italic,
  BoldItalicFont = *-bolditalic,
]
\setmonofont{texgyrecursor}[
  Extension      = .otf,
  UprightFont    = *-regular,
  BoldFont       = *-bold,
  ItalicFont     = *-italic,
  BoldItalicFont = *-bolditalic,
]

% Vietnamese, imported from babel's own locale data and made the main language.
% This must come after fontspec so the language activates against fonts that
% are already set.
\babelprovide[import, main]{vietnamese}
